\section{The finite-interval action inequality}\label{sec:action-framework}

We first formulate the lower bound on the variational action that will
lead to an upper bound on spectral moments. Let $\Omega=(x_-,x_+)$ be a bounded interval,
let $n,d\ge1$, and suppose
\begin{equation}\label{eq:row-orthonormality}
 U\in H^1_0(\Omega;\R^{n\times d}),\qquad
 \int_\Omega U(x)U(x)^{\mathsf T}\,dx=I_n.
\end{equation}
Thus the rows of $U$ are orthonormal as $\R^d$-valued functions. Fix
$K=\diag(k_1,\ldots,k_n)$ with each $k_i>0$. For
$A\in H^1_0(\Omega;\R^{n\times d})$ define
\begin{equation}\label{eq:action}
 \mathcal E(A)=\int_\Omega
 \left(\|A'\|_{\HS}^2+\tr(K^2AA^{\mathsf T})
                  -\tr((AA^{\mathsf T})^q)\right)\,dx.
\end{equation}
This integral is finite, since one-dimensional $H^1$ functions on a bounded
interval are continuous. Write $\mathcal L$ for the real lower triangular
$n\times n$ matrices.

\begin{theorem}[Action inequality]\label{thm:action}
For every $0<\sigma<1$ and all the data just specified,
\begin{equation}\label{eq:action-bound}
 \sup_{C\in\mathcal L}\mathcal E(CU)
 \ge\sum_{i=1}^n\int_{-k_i}^{k_i}(k_i^2-s^2)^\sigma\,ds.
\end{equation}
\end{theorem}

The order of the $k_i$ is immaterial for this theorem. For its spectral
application we order them decreasingly: lower triangularity then preserves
the nested trial spaces on which the form is at most $-k_i^2$ times the
squared norm. Matrix Young duality bounds the action from above by a constant
times $\int\tr(W^p)$, while
\[
 \int_{-k_i}^{k_i}(k_i^2-s^2)^\sigma\,ds
 =\mathrm B(\tfrac12,1+\sigma)k_i^{2\sigma+1}
\]
has exactly the desired spectral exponent. Section~\ref{sec:spectral}
gives this deduction, including its approximation steps.

The scalar companion proves the version with $d=1$
\cite[Theorem~2.1]{ScalarCompanion}. Its field and endpoint construction
also suggest the present proof. The necessary new ingredient is an
inequality for positive matrix densities of unrestricted rank. We prove
that comparison next, before constructing the regularized field to which
it applies. The proof of Theorem~\ref{thm:action} is completed in
Section~\ref{sec:action-limit}.
